% Part: first-order-logic
% Chapter: completeness
% Section: introduction

\documentclass[../../../include/open-logic-section]{subfiles}

\begin{document}

\iftag{FOL}
      {\olfileid{fol}{com}{int}}
      {\olfileid{pl}{com}{int}}

\olsection{పరిచయం}

సంపూర్ణతా సిద్ధాంతం తర్కం గురించిన అత్యంత మౌలిక ఫలితాలలో ఒకటి.
దానికి రెండు సూత్రీకరణలు ఉన్నాయి; వాటి సమానార్థకతను మనం
నిరూపిస్తాం. మొదటి సూత్రీకరణలో అది అర్థపర అనుగమనం, మన
!!{derivation} వ్యవస్థల మధ్య సంబంధం గురించి ఒక మౌలిక విషయాన్ని
చెబుతుంది: కొన్ని !!{sentence}s~$\Gamma$ నుండి !!a{sentence}~$!A$
అనుగమిస్తే, $\Gamma \Proves !A$ను స్థాపించే !!a{derivation} కూడా
ఉంటుంది. ఆ విధంగా, వాస్తవంగా అనుగమించని విషయాలను నిరూపించకుండానే
సాధ్యమైనంత బలంగా మన !!{derivation} వ్యవస్థ ఉంటుంది.

రెండవ సూత్రీకరణలో దాన్ని ఒక నమూనా-ఉనికి ఫలితంగా చెప్పవచ్చు: ప్రతి
అవైరుధ్యమైన !!{sentence}s సమితి సంతృప్తిపరచదగినది. అవైరుధ్యం ఒక
నిరూపణ-సిద్ధాంత భావన: మన !!{derivation} వ్యవస్థ కొన్ని రకాల
!!{derivation}sను ఉత్పత్తి చేయలేదని అది చెబుతుంది. అయితే~$\Gamma$
నుండి ఒక నిర్దిష్ట రకమైన !!{derivation}s లేవనే కారణంతోనే
$\iftag{FOL}{\Sat{M}}{\pSat{v}}{\Gamma}$ అయ్యే
\iftag{FOL}{!!a{structure}~$\Struct{M}$}{!!{valuation}{~$\pAssign{v}$}}
ఉందని ఎవరు హామీ ఇవ్వగలరు? సంపూర్ణతా సిద్ధాంతం మొదటిసారిగా
నిరూపించబడకముందు---నిజానికి ఇప్పుడు మనకు ఉన్న !!{derivation}
వ్యవస్థలు ఏర్పడకముందే---గొప్ప జర్మన్ గణితశాస్త్రవేత్త డేవిడ్
హిల్బర్ట్, గణిత సిద్ధాంతాల అవైరుధ్యం అవి చెప్పే వస్తువుల ఉనికికి
హామీ ఇస్తుందని భావించాడు. గాట్లోబ్ ఫ్రేగెకు రాసిన ఒక లేఖలో ఆయన
ఇలా చెప్పాడు:
\begin{quote}
  యాదృచ్ఛికంగా ఇచ్చిన స్వీకృతాలు, వాటి పరిణామాలన్నింటితో సహా,
  పరస్పరం వైరుధ్యపడకపోతే అవి సత్యమైనవి; ఆ స్వీకృతాలు నిర్వచించే
  వస్తువులు ఉనికిలో ఉంటాయి. ఇదే నాకు సత్యం, ఉనికికి ప్రమాణం.
\end{quote}
ఫ్రేగె దీనిని తీవ్రంగా విభేదించాడు. స్వీకృతాలు అవైరుధ్యంగా ఉంటే,
వాటన్నింటినీ సత్యం చేసే \emph{ఏదో ఒక}
\iftag{FOL}{!!{structure}}{!!{valuation}} ఉంటుందనే అర్థంలోనైనా
హిల్బర్ట్ సరైనవాడని సంపూర్ణతా సిద్ధాంతపు రెండవ సూత్రీకరణ చూపుతుంది.

సంపూర్ణతా సిద్ధాంతం---లేదా ఇంకా సరిగ్గా చెప్పాలంటే, దాని
నిరూపణ---ముఖ్యమైనది కావడానికి ఇవి మాత్రమే కారణాలు కావు. దానికి అనేక
ముఖ్యమైన పరిణామాలు ఉన్నాయి; వాటిలో కొన్నింటిని వేరుగా చర్చిస్తాం.
ఉదాహరణకు, $\Gamma \Proves !A$ను చూపే ఏ !!{derivation} అయినా పరిమితమే;
అందువల్ల అది~$\Gamma$లోని !!{sentence}sలో పరిమిత సంఖ్యలో ఉన్నవాటినే
ఉపయోగించగలదు. కనుక~$!A$ అనేది~$\Gamma$ యొక్క పరిణామమైతే, అది ఇప్పటికే
$\Gamma$లోని ఒక పరిమిత ఉపసమితి యొక్క పరిణామమని సంపూర్ణతా సిద్ధాంతం
ద్వారా అనుగమిస్తుంది. దీనిని \emph{సంహతత్వం} అంటారు. సమానార్థకంగా,
$\Gamma$లోని ప్రతి పరిమిత ఉపసమితి అవైరుధ్యమైతే, $\Gamma$ కూడా
అవైరుధ్యంగానే ఉండాలి.

సంహతత్వ సిద్ధాంతం సంపూర్ణతా సిద్ధాంతం నుంచి !!{derivation}s మీదుగా
వెళ్లే మార్గంలో అనుగమించినప్పటికీ, దాన్ని నేరుగా స్థాపించడానికి
సంపూర్ణతా సిద్ధాంతపు \emph{నిరూపణనే} ఉపయోగించవచ్చు. ఎందుకంటే ఆ నిరూపణ
ఒక నిర్దిష్ట ధర్మం---అవైరుధ్యం---గల !!{sentence}s సమితిని తీసుకొని,
దానిలో నుంచి కొన్ని ధర్మాలు గల !!a{structure}ను నిర్మిస్తుంది (ఈ
సందర్భంలో, ఆ సమితిని అది సంతృప్తిపరుస్తుంది). అవైరుధ్య సమితులకు బదులుగా
``పరిమితంగా సంతృప్తిపరచదగిన'' !!{sentence}s సమితులతో మొదలుపెడితే,
దాదాపు అదే నిర్మాణాన్ని ఉపయోగించి సంహతత్వాన్ని నేరుగా స్థాపించవచ్చు.
\iftag{FOL}{ఈ నిర్మాణం మరికొన్ని పరిణామాలను కూడా ఇస్తుంది; ఉదాహరణకు,
  సంతృప్తిపరచదగిన ఏ !!{sentence}s సమితికైనా పరిమిత లేదా
  !!{denumerable}s నమూనా ఉంటుందని చూపుతుంది. (ఈ ఫలితాన్ని
  లొవెన్‌హైమ్--స్కోలెమ్ సిద్ధాంతం అంటారు.) సాధారణంగా, !!{sentence}s
  సమితుల నుంచి !!{structure}sను నిర్మించే పద్ధతిని తర్కంలో తరచుగా,
  కొన్నిసార్లు తత్వశాస్త్రంలో కూడా, ఉపయోగిస్తారు.}{}

\end{document}
